% Zone „Das kennst du schon“ – aus bank/unabhaengigkeit/zone.jsonl (zusammenbau v0.1)
\einheitenkopf[zone][Kennst du schon]{Das kennst du schon}
\setcounter{aufgabe}{0}

%% TODO Titel: Ich-kann-Satz fehlt in der Bank (Platzhalter: Kettenname) – Nr. 1
%% TODO Anweisung: ein Satz über der ersten Teilaufgabe fehlt in der Bank – Nr. 1
\begin{aufgabe}{Vierfeldertafel füllen (Ränder, Felder, bedingte Angaben)}
\begin{teile}
\teil Von 40 Schülern sind 24 Mädchen. 10 Schüler tragen eine Brille, darunter 6 Mädchen. Ergänze die Tafel.

\vierfeldertafel{Mädchen}{Brille}{6,,10,,,,24,,40}
\teil In einer Umfrage sind 55\,\% der Befragten weiblich. 30\,\% aller Befragten fahren regelmäßig Rad, 18\,\% sind weiblich und fahren Rad. Ergänze die Tafel mit Prozentwerten.

\vierfeldertafel{Weiblich}{Rad}{18\,\%,,30\,\%,,,,55\,\%,,100\,\%}
\end{teile}
\end{aufgabe}

%% TODO Titel: Ich-kann-Satz fehlt in der Bank (Platzhalter: Kettenname) – Nr. 2
%% TODO Anweisung: ein Satz über der ersten Teilaufgabe fehlt in der Bank – Nr. 2
\begin{aufgabe}{Bedingte Wahrscheinlichkeit als Quotient bilden}
\begin{teile}
\teil Von 20 Schülern haben 8 ein Fahrrad, darunter 5 Mädchen. Aus den Schülern mit Fahrrad wird einer ausgelost. Mit welcher Wahrscheinlichkeit ist es ein Mädchen? \leerfeld
\teil Für zwei Ereignisse gilt $P(A) = 0{,}4$ und $P(A \cap B) = 0{,}1$. Berechne $P_A(B)$. \leerfeld
\end{teile}
\end{aufgabe}

%% TODO Titel: Ich-kann-Satz fehlt in der Bank (Platzhalter: Kettenname) – Nr. 3
%% TODO Anweisung: ein Satz über der ersten Teilaufgabe fehlt in der Bank – Nr. 3
\begin{aufgabe}{Anteile aus absoluten Häufigkeiten bilden und kürzen}
\begin{teile}
\teil 12 von 48 Schülern – welcher Anteil? \leerfeld
\teil Von 1\,250 Prüflingen haben 875 bestanden. Gib den Anteil als Dezimalzahl an. \leerfeld
\end{teile}
\end{aufgabe}

%% TODO Titel: Ich-kann-Satz fehlt in der Bank (Platzhalter: Kettenname) – Nr. 4
%% TODO Anweisung: ein Satz über der ersten Teilaufgabe fehlt in der Bank – Nr. 4
\begin{aufgabe}{Pfadterme mit Parameter im Baumdiagramm (Ast als Faktor, Pfad als Produkt, Randwahrscheinlichkeit als Pfadsumme)}
\begin{teile}
\teil Ein Glücksrad zeigt Rot mit der Wahrscheinlichkeit p, sonst Blau. Es wird zweimal gedreht. Gib den Term für die Wahrscheinlichkeit von „zweimal Rot“ an.
\teil Ein Glücksrad zeigt Rot mit der Wahrscheinlichkeit p, sonst Blau; es wird zweimal gedreht. Welcher Term gehört zu „genau einmal Rot“? \\ \kreuz{$2 \cdot p \cdot (1 - p)$} \\ \kreuz{$p \cdot (1 - p)$} \\ \kreuz{$2p$}
\end{teile}
\end{aufgabe}

%% TODO Titel: Ich-kann-Satz fehlt in der Bank (Platzhalter: Kettenname) – Nr. 5
%% TODO Anweisung: ein Satz über der ersten Teilaufgabe fehlt in der Bank – Nr. 5
\begin{aufgabe}{Lineare und quadratische Gleichungen aufstellen und lösen, Lösungen sieben (Nulllösung verwerfen)}
\begin{gleichungsraster}[2]
\gl{\text{Löse: $0{,}3x = 0{,}12$}} &
\gl{\text{Löse: $x \cdot (0{,}5 - x) = 0{,}06$}} \\
\end{gleichungsraster}
\end{aufgabe}

%% TODO Titel: Ich-kann-Satz fehlt in der Bank (Platzhalter: Kettenname) – Nr. 6
%% TODO Anweisung: ein Satz über der ersten Teilaufgabe fehlt in der Bank – Nr. 6
\begin{aufgabe}{Pfadregeln und Gegenereignis mit und ohne Zurücklegen}
\begin{teile}
\teil Ein Würfel wird zweimal geworfen. Mit welcher Wahrscheinlichkeit fällt zweimal eine Sechs? \leerfeld
\teil Von 5 Losen sind 2 Gewinne. Zwei Lose werden nacheinander ohne Zurücklegen gezogen. Mit welcher Wahrscheinlichkeit sind beide Gewinne? \leerfeld
\end{teile}
\end{aufgabe}

%% TODO Titel: Ich-kann-Satz fehlt in der Bank (Platzhalter: Kettenname) – Nr. 7
%% TODO Anweisung: ein Satz über der ersten Teilaufgabe fehlt in der Bank – Nr. 7
\begin{aufgabe}{Vierfeldertafel füllen (Ränder, Felder, bedingte Angaben) – Fehler finden}
\begin{teile}
\teil Von 50 Mitgliedern eines Vereins sind 20 Jugendliche. 30\,\% der Jugendlichen spielen Tennis; insgesamt spielen 15 Mitglieder Tennis. Für das Feld „Jugendlich und Tennis“ rechnet Ole: \rechnung{0{,}3 \cdot 50 &= 15} Finde den Fehler.
\end{teile}
\end{aufgabe}

%% TODO Titel: Ich-kann-Satz fehlt in der Bank (Platzhalter: Kettenname) – Nr. 8
%% TODO Anweisung: ein Satz über der ersten Teilaufgabe fehlt in der Bank – Nr. 8
\begin{aufgabe}{Vierfeldertafel füllen (Ränder, Felder, bedingte Angaben) – selbst rechnen}
\begin{teile}
\teil Von 80 Besuchern eines Kinos sind 32 Kinder. 25\,\% der Kinder kaufen Popcorn; insgesamt kaufen 30 Besucher Popcorn. Wie viele Erwachsene kaufen Popcorn? \leerfeld Erwachsene
\end{teile}
\end{aufgabe}
